\begin{abstract}
CuTe, the core of NVIDIA's CUTLASS, describes how tensors, threads and hardware instructions index
memory by \emph{layouts}---nested (shape, stride) pairs---and manipulates them with an algebra of
partial operations: coalesce, composition, complement, logical division and product. The operations
are specified by post-conditions and guarded by divisibility checks, and existing accounts give
sufficient conditions for them to be defined. We show that the algebra is governed by a single
identity. Writing a layout function in ``slope--jump'' form,
$f(x)=x\,d_0+\sum_i\lfloor x/P_i\rfloor\,\gamma_i$, one has
$f(x+y)=f(x)+f(y)+\sum_i c_i(x,y)\,\gamma_i$, where $c_i$ is the carry into digit $i$ of the
mixed-radix system given by the shape: a layout is additive up to carries, and each carry is
charged a jump $\gamma_i$. From this we obtain (1) a characterisation of layout functions---the
M\"obius transform of their first differences is supported on a divisibility chain---which yields
the coalesced normal form and a recognition algorithm; (2) an exact criterion for composition:
the composite, when it exists, is obtained by evaluating the outer layout at the strides of the
inner one, and it exists if and only if a ``carry defect'' vanishes; CuTe's stride and shape
divisibility conditions describe a special case, that of strides with a single non-zero digit
dividing the size of its mode; (3) the identification of complements with the classical
direct-sum decompositions of an interval (Krasner, de~Bruijn, Long), and an exact criterion for
logical division. The results on normal forms and composition hold for strides in any abelian
group,
covering negative strides, coordinate-valued layouts and XOR-based (Triton-style) layouts. Applied
to the binary radix, the same identity covers CuTe's swizzles: a swizzle is a layout for the XOR
group law on offsets; we characterise when a swizzled layout without offset is XOR-linear, give a
closed-form
criterion, under a bit-disjointness assumption, for a slice of it to decay to an ordinary layout,
and give conditions under which CuTe's operations on swizzled layouts are valid. The criteria explain two behaviours of current implementations that we observe in differential
testing: composites that are silently wrong because modes of the inner layout interfere, and
composable pairs that are rejected. The carry identity, the characterisation of layout functions,
the composition formula with its defect and the carry-free criterion for the naive composite are
machine-checked in Lean~4; the remaining results have paper proofs and are validated by exhaustive and randomised
checks.
\end{abstract}
